%% ma: L10-B13-0056
%% lop: 10
%% chuong: 5
%% bai: 13
%% dang: 10.13.2
%% loai: TLN
%% muc: TH
%% dap_an: "6,23"
%% nguon: "chuyen-de-cac-so-dac-trung-cua-mau-so-lieu-khong-ghep-nhom-toan-10-chuong-trinh-moi.pdf (GV. Phan Nhật Linh) (P02), Bài 02 Dạng 1 Phần III Câu 6"
%% kien_thuc: "Lớp 10 Bài 13 (số trung bình)"
%% trang_thai: chinh-thuc
%% ghi_chu: "Tự duyệt theo yêu cầu giáo viên 22/10/2026."
\begin{ex}
Để khảo sát kết quả môn Toán (thang điểm 10) trong kì thi tuyển sinh năm vừa qua của trường A, người ta chọn một mẫu gồm $100$ học sinh. Điểm môn Toán của các học sinh được cho ở bảng tần số sau:
\begin{center}
\begin{tabular}{|c|c|c|c|c|c|c|c|c|c|c|c|}\hline
Điểm & $0$ & $1$ & $2$ & $3$ & $4$ & $5$ & $6$ & $7$ & $8$ & $9$ & $10$\\\hline
Tần số & $1$ & $1$ & $3$ & $5$ & $8$ & $13$ & $19$ & $24$ & $14$ & $10$ & $2$\\\hline
\end{tabular}
\end{center}
Số trung bình của mẫu số liệu trên là bao nhiêu?
\shortans{6,23}
\loigiai{$\overline x=\dfrac{0\cdot1+1\cdot1+2\cdot3+\cdots+10\cdot2}{100}=\dfrac{623}{100}=6{,}23$.}
\end{ex}
